\PassOptionsToPackage{table, dvipsnames}{xcolor}
\usepackage{xcolor}
\usepackage[utf8]{inputenc}
\usepackage[T1]{fontenc}
\usepackage[brazilian]{babel}

\usepackage{mathptmx}
\usepackage{setspace}
\usepackage{indentfirst}
\onehalfspacing

\usepackage{geometry}
\geometry{
    a4paper,
    left=3cm,
    top=3cm,
    right=2cm,
    bottom=2cm
}

\usepackage{amsmath}
\usepackage{amsfonts}
\usepackage{amssymb}
\usepackage{siunitx}
\sisetup{
    output-decimal-marker = {,},
    per-mode = symbol,
    separate-uncertainty = true
}

\usepackage{graphicx}
\usepackage{float}
\usepackage{caption}
\usepackage{subcaption}
\usepackage{booktabs}
\usepackage{array}
\usepackage{multirow}
\usepackage{longtable}
\usepackage{xurl}
\usepackage[
    colorlinks=true,
    linkcolor=black,
    citecolor=black,
    urlcolor=black
]{hyperref}

\graphicspath{{./pags/imagens/}}

% ------------------------------------------------------------
% MARCADOR DE DADOS A PREENCHER (aparece em vermelho no PDF)
% Funciona em texto e em modo matemático. Ex.: \ph{V_1}, \ph{\square}
% Quando tiver os dados, troque \ph{...} pelo valor numérico.
% Para achar tudo que falta, busque por "\ph" no editor.
% Para entregar sem vermelho, troque a definição por:
%   \newcommand{\ph}[1]{\ensuremath{#1}}
% ------------------------------------------------------------
\newcommand{\ph}[1]{\textcolor{red}{\ensuremath{#1}}} % marcador de dados a preencher (não usado na versão final)
